%
%============= CHAPTER BREAK =================%
%
\chapter{Antigonus of Nicaea (ca 150 CE)}%
\vspace*{3mm}%
\justifying%
%
\noindent{}Antigonus of Nicaea, or Antigonus the Astrologer, %
wrote his \textit{Horoscope of Hadrian} / \textit{Hadriani genitura} around 150 CE.\endnote{%
Heilan 2015:1.} %
This work happens to be the earliest extant biographical source on the life of the emperor Hadrian.\endnote{%
Heilan 2015:2.} \\[1mm] %

\noindent{}\textbf{κοσμοκράτωρ} / \textit{kosmokratōr} / “world-ruler” is a term well known in the history of Christian %
dogma and liturgy, but did Christians invent the term or appropriate it from their predecessors? %
A TLG query of the lemma shows that Antigonus and Ptolemy attest its %
earliest extant non-Christian uses.\endnote{%
A fragment attributed to Ctesias (5th--4th cent. BC) is the topmost hit in the %
date-sorted list of results, but it is used within a phrase with overt Christian language, %
and is thus not a viable attestation or candidate for the lemmaʼs point of origin.} %
The relevant passage in Antigonus follows immediately after his report of the death of Hadrian %
from “dropsy” / ὑδρωπικὸς.\endnote{%
Kroll 6:68.} This famous death, incidentally, %
as well as its surrounding astrological language, %
may have partly inspired the Lukan story of Jesus healing “a certain man with dropsy” / %
ἄνθρωπός τις ἦν ὑδρωπικὸς (Lk 14:2), %
whom he encountered “in the house of a certain one of the rulers” / %
εἰς οἶκόν τινος τῶν ἀρχόντων (Lk 14:1). %
Antigonus proceeds to describe the astrologically auspicious birth of Hadrian:\endnote{%
Kroll 6:68.} %

\vspace*{2mm} 
 \begin{quote}καὶ αὐτῶν τῶν δορυφορούντων αὐτῇ ἀστέρων οἰκείως εὑρεθέντων, %
 τῆς μὲν Ἀφροδίτης ἐν οἰκείῳ ὑψώματι, %
 τοῦ δὲ Ἄρεως ἐν ἰδίῳ τριγώνῳ καὶ ἰδίαις μοίραις τὴν στάσιν ἔχοντος, %
 ἀμφοτέρων ἰδιοτοπούντων καὶ ἐπαναφερομένων τῇ Σελήνῃ· %
 προσέτι δὲ καὶ τοῦ \textbf{κοσμοκράτορος} Ἡλίου δορυφοροῦντος αὐτῇ ἐν ταῖς ἑξῆς αὐτῆς %
 μοίραις καὶ αὐτοῦ δορυφορουμένου παρά τε τοῦ Κρόνου ὄντος ἐν τῷ ἰδίῳ οἴκῳ καὶ Ἑρμοῦ, %
 ἀμφοτέρων ὄντων ἐπὶ ἑῴας ἀνατολῆς. \\[1mm] %
 And of the stars who guard-attend her that are familiarly found, %
 that of Aphrodite in familiar exaltation, %
 that of Ares in his proper triangle and holding place with his proper Fates, %
 both of proper place and counter-revolving with the moon. %
 Over and above, also that of the \textbf{world-ruler} Helios guard-attending her %
 in her successive Fates, and he himself attended by Cronos in his proper house and Hermes, %
 both of whom [are] on Eastʼs horizon. %
 \end{quote} % 
 \vspace*{2mm}

\noindent{}Antigonusʼ use of the term in the singular makes it distinct from that of Ptolemy, %
who uses the plural instead, specifically in a section on %
“worth and ... happiness” / ἀξίας καὶ ... εὐδαιμονίας: %

\vspace*{2mm} 
 \begin{quote}καὶ ἐὰν μὲν οἱ δορυφοροῦντες ἀστέρες ἤτοι ἐπίκεντροι καὶ %
 αὐτοὶ ὦσιν ἢ πρὸς τὸ ὑπὲρ γῆν κέντρον συσχηματίζωνται, %
 μεγάλοι καὶ δυναμικοὶ καὶ \textbf{κοσμοκράτορες} διατελοῦσι, %
 καὶ ἔτι μᾶλλον εὐδαίμονες, %
 ἐὰν οἱ δορυφοροῦντες δεξιοὶ τοῖς ὑπὲρ γῆν κέντροις συσχηματίζωνται, %
 ἐὰν δὲ τῶν ἄλλων οὕτως ἐχόντων μόνος ὁ ἥλιος ᾖ ἐν ἀρρενικῷ ἡ δὲ σελήνη ἐν θηλυκῷ, %
 ἐπίκεντρον δὲ τὸ ἕτερον τῶν φώτων, %
 ἡγεμόνες μόνον ἔσονται ζωῆς καὶ θανάτου κύριοι. (\textit{Tetr}. 4.3.2) \\[1mm] %
 \textls[-10]{And if the guard-attending stars either be themselves at cardinal points %
 or are arranged together unto a point over earth, %
 great and powerful and \textbf{world-rulers} are brought forth, %
 and even more fortunate, %
 if the northerly guard-attending are arranged together by the points over earth. %
 But of all others who have this, only if the sun be in a masculine position, %
 and the moon in a feminine position, %
 and another cardinal point of the stars, %
 only then will rulers be lords of life and death.} %
 \end{quote} % 
\vspace*{2mm}

\noindent{}Despite the numerical difference in the use of this keyword, %
the passages overlap considerably in regard to subject (auspicious births of rulers), %
technical astrological terms (“guard-attending”), %
and specific stars. %

The same term appears in a variety of 2nd century Christian texts, %
and all of them are most likely subsequent to and perhaps even dependent upon %
(even if loosely and/or combatively) Antigonus and/or Ptolemy. %
Both Marcionite Laodiceans and canonical Ephesians use %
the term in the plural, closest to Ptolemy in this regard, %
and both have surrounding terminology highly evocative of astrological imagery. %
Word differences are rendered in italics to facilitate comparison (Lao 6:12): %

\vspace*{3mm} 
 \begin{quote}\textls[-15]{ὅτι ἔστιν \textit{ὑμῖν} ἡ πάλη %
  πρὸς τὰς ἐξουσίας, πρὸς τοὺς \textbf{κοσμοκράτορας} τοῦ σκότους τούτου, %
  πρὸς τὰ πνευματικὰ τῆς πονηρίας ἐν τοῖς ἐπουρανίοις.} \\[1mm] % %
  Because for \textit{you} the battle is before the authorities, %
  before the \textbf{world-rulers} of this darkness, %
  before the spiritualia of evil in the heavenlies. %
  \end{quote} % 
 \vspace*{1.5mm}

 \begin{quote}ὅτι \textit{οὐκ} ἔστιν ἡμῖν ἡ πάλη \textit{πρὸς αἷμα καὶ σάρκα ἀλλὰ πρὸς τὰς ἀρχάς}, %
  πρὸς τὰς ἐξουσίας, πρὸς τοὺς \textbf{κοσμοκράτορας} τοῦ σκότους τούτου, %
  πρὸς τὰ πνευματικὰ τῆς πονηρίας ἐν τοῖς ἐπουρανίοις. (Eph 6:12) \\[1mm] %
  Because for \textit{us} the battle is \textit{not before blood and flesh, %
  but instead before the rulers}, %
  before the authorities, before the \textbf{world-rulers} of this darkness, %
  before the spiritualia of evil in the heavenlies. %
  \end{quote} % 
 \vspace*{3mm}

Though it is derivative, a quotation of Eph 6:12 in 13Ign \textit{Eph} 11:2 also attests this plural form. %
The term appears in the singular form in Ps-Clement, \textit{Hom}. 6.21.2 %
and \textit{ActJohn} 23. %